\clearpage
\section{Count-based language models}\label{sec:counts}
\subsection*{Approximate the context, then estimate probabilities}
The chain rule permits the full prefix. An $n$-gram approximation keeps
only the last $n-1$ tokens:
\[
 p(x_t\mid x_{<t})\approx p(x_t\mid x_{t-n+1:t-1}).
\]
A unigram has no token context, a bigram has one context token, and a
trigram has two. ``Trigram'' counts the predicted token as well as its context.
The approximation concerns conditional independence, not the factorization itself.

\fig{chain-context}{A prediction for the final token can use increasingly long
histories. The full-prefix chain rule remains exact; the short histories impose a model assumption.}

For a history $h$, let $C(h,w)$ count training occurrences followed by $w$,
and let $C(h)=\sum_{v\in\mathcal V}C(h,v)$. Maximum likelihood gives
\[
 \widehat p(w\mid h)=\frac{C(h,w)}{C(h)},\qquad C(h)>0.
\]
To see why, minimize $-\sum_w C(h,w)\log p_w$ subject to $\sum_wp_w=1$.
The Lagrange multiplier condition is $-C(h,w)/p_w+\lambda=0$ for observed
events; normalization gives $\lambda=C(h)$.

\fig{count-probability}{The course example has three occurrences of history
\texttt{I}: two followed by \texttt{am}, one by \texttt{do}. Counts and probabilities have different units.}

We extract adjacent pairs separately inside each sentence, including the
chosen boundary markers. Joining sentences first could invent a pair from
the last word of one sentence to the first word of the next. In the
L02/L03 example, \BOS{} is context-only and \EOS{} is a possible target.
The row for \texttt{I} assigns $2/3$ to \texttt{am} and $1/3$ to
\texttt{do}; all other targets receive zero under unsmoothed MLE.

\selfcheck{Which count belongs in the denominator: occurrences of the predicted
word or occurrences of the history? Explain why each conditional row must
sum to one, even if a plotted table puts histories in columns instead.}

\nextpage{Sparse counts, smoothing, and unseen histories}
A zero training count is often weak evidence that an event is impossible.
Yet unsmoothed MLE assigns zero probability and infinite loss to such an
event. Additive smoothing reserves probability for every symbol in a
fixed output vocabulary:
\[
 p_\delta(w\mid h)=\frac{C(h,w)+\delta}{C(h)+\delta K},\qquad\delta>0.
\]
The denominator is the sum of the numerators, so normalization is explicit.
For an unseen history, $C(h)=0$ gives the uniform distribution $1/K$.
This is a defined fallback, unlike the undefined MLE ratio $0/0$.

\fig{smoothing}{An original three-symbol example with counts $(2,1,0)$.
Add-one smoothing changes the row to $(1/2,1/3,1/6)$; interpolation uses a
chosen uniform fallback with weight $0.3$.}

For a mixture of normalized component models,
\[
 p(w\mid h)=\sum_{r=0}^{n}\lambda_r p_r(w\mid h),\qquad
 \lambda_r\ge0,\quad\sum_r\lambda_r=1,
\]
the result is normalized because summation over $w$ commutes with the finite
mixture. Components can use different history lengths; $p_0$ can be a
uniform floor. Each component must define its behavior for unseen histories.
Tune mixture weights on development text, then hold them fixed for the test set.

\textbf{A useful distinction: frequency versus continuation diversity.}
Imagine a word that appears 100 times, always after the same preceding
word, and another that appears only ten times after ten different words.
The latter has appeared in more kinds of contexts. Kneser--Ney smoothing
uses continuation statistics for its lower-order distribution, rather than
simply recycling unigram frequency. \readingref{R02} develops this idea;
the recap requires the distinction, not a full implementation.

\textbf{Unknown words and output support.} A word tokenizer may map unseen
forms to \UNK{}. A byte tokenizer can represent their bytes, but the language
model can still encounter unseen byte or subword contexts. Input-only
markers should not silently enter the output denominator. Vocabulary size
means the number of outcomes the model actually predicts in this example.

\takeaway{A language model needs a complete normalized rule for every
history encountered at evaluation or generation time.}
\selfcheck{For $\delta=1$ and counts $(2,1,0)$, verify the unseen symbol gets
$1/6$. What happens as $\delta\to\infty$? The row approaches the uniform distribution.}
